% ==============================================================================
% =================================== Aufgabe 1 =================================
% ==============================================================================

\chapter{Aufgabe}
\label{chap:1 Aufgabe}
\thispagestyle{fancy}
\section{Was zeichnet zukunftsfähige ERP-Systeme aus? Nennen Sie fünf Erfolgsfaktoren hierfür.} \par \uline{Zukunftsfähige \acs{ERP}-Systeme} zeichnen sich primär durch eine \uline{hohe Wandlungsfähigkeit} aus, also die Fähigkeit, sich effizient und schnell an veränderte Anforderungen anzupassen. \cite{1}\cite{2} In einem entsprechenden Bewertungsportfolio liegt der Idealbereich für zukunftsfähige Systeme dort, wo sowohl eine \uline{hohe technische} als auch eine \uline{hohe geschäftsspezifische Wandlungsfähigkeit} vorliegt. \cite{3}\cite{4} \par Fünf entscheidende Erfolgsfaktoren für solche Systeme sind:
\begin{enumerate}[label=\arabic*.), leftmargin=*]
    \item \textbf{Hohe Wandlungsfähigkeit (\textit{Adaptivität}):} Ein zukunftsfähiges System muss in der Lage sein, sich ohne kompletten Austausch an neue Marktanforderungen und organisatorische Veränderungen anzupassen. \cite{5}\cite{6} Dies wird technisch durch Kriterien wie \uline{Skalierbarkeit}, \uline{Modularität} und \uline{Interoperabilität} sowie durch autopoietische Eigenschaften wie \uline{Selbstorganisation} unterstützt. \cite{7}\cite{8}\cite{9}\cite{10}
    \item \textbf{Integration von Künstlicher Intelligenz (\textit{\acs{KI}}) und Analytics:} Zukunftsfähige Systeme gehen über eine rein vergangenheitsorientierte Sicht hinaus und integrieren \uline{Monitoring-}, \uline{Prognose-}, \uline{Simulations-} und \uline{Optimierungsfunktionen}. \cite{11}\cite{12} Die Nutzung von \acs{KI} ermöglicht eine unvoreingenommene Analyse vorhandener Daten, bietet verlässliche Vorhersagen und unterstützt die Entscheidungsfindung durch digitale Assistenten. \cite{12}\cite{13}\cite{14}
    \item \textbf{Rekonfigurierbare Geschäftsprozesse (\textit{"Modellieren statt Programmieren"}):} Aufgrund der zunehmenden \uline{Delinearisierung} von Arbeitsabläufen müssen zukünftige \acs{ERP}-Systeme Prozessmodelle aufweisen, die von den Nutzern zur Laufzeit individuell angepasst werden können. \cite{15}\cite{16} Die Einbindung von \uline{Workflow-Designern} ermöglicht es, Standardprozesse flexibel um spezialisierte Abläufe zu ergänzen. \cite{16}
    \item \textbf{Interoperabilität und Delokalisierung:} In einer vernetzten Welt müssen \acs{ERP}-Systeme die Unternehmensgrenzen öffnen und den Datenaustausch über die gesamte Wertschöpfungskette hinweg ermöglichen. \cite{11}\cite{17} Die Schnittstellen müssen \uline{interoperabel} und \uline{selbstbeschreibungsfähig} sein, um eine nahtlose Kommunikation zwischen unterschiedlichen Objekten des digitalen Unternehmens zu gewährleisten. \cite{11}\cite{18}
    \item \textbf{Skalierbare Infrastruktur und Cloud-Fähigkeit:} Ein zukunftsfähiges System muss die \uline{Skalierbarkeit des Geschäftsbetriebs} sowohl langfristig (\textit{Wachstum}) als auch kurzfristig (\textit{Saisongeschäft}) unterstützen. \cite{19} Moderne Cloud-Architekturen und hybride Datenumgebungen sind dabei essenziell, um die erforderliche Leistung für transaktionale und analytische Prozesse bereitzustellen. \cite{20}\cite{21}
\end{enumerate}
\section{Entwerfen Sie ein Projektteam für die innerbetriebliche Migration des \acs{ERP}-Systems. Beschreiben Sie die wichtigsten Ziele des Projekts.} \par Für die innerbetriebliche Migration (\textit{Einführung}) eines \acs{ERP}-Systems setzt sich ein effektives \uline{Projektteam} aus verschiedenen hierarchischen und funktionalen Rollen zusammen:
\begin{itemize}
    \item Der \uline{Lenkungsausschuss} bildet das oberste Entscheidungsgremium und besteht idealerweise aus der Geschäftsführung, dem Gesamtprojektleiter, dem Projektleiter des Anbieters sowie einem \uline{Trusted Advisor}. \cite{22}
    \item Die Aufgaben dieses Ausschusses umfassen die Festlegung und Veränderung von Projektzielen, Terminen und Budgets sowie die Harmonisierung des Projekts mit den strategischen Unternehmenszielen. \cite{23}
    \item Ein \uline{Trusted Advisor} fungiert dabei als Ratgeber der Geschäftsführung, bewertet die Qualität der Ergebnisse und stellt sicher, dass Anbieter und Kunde eine gemeinsame Sprache sprechen. \cite{22}\cite{24}
    \item Der \uline{Gesamtprojektleiter} trägt die operative Verantwortung, berichtet an den Lenkungsausschuss, genehmigt Projektphasen und überwacht die Kompatibilität zwischen der \acs{ERP}-Lösung und den Organisationsanforderungen. \cite{25}
    \item \uline{Key User} (\textit{Hauptanwender}) wirken als Multiplikatoren in ihren jeweiligen Fachbereichen, definieren Anforderungen in Workshops und leisten eine aktive Unterstützung bei der Konzeptumsetzung. \cite{25}
    \item Die eigentlichen \uline{Projektmitarbeiter} sollten zu mehr als 50\% ihrer Arbeitszeit für das Projekt freigestellt sein, da sie die Verantwortung für die Konzepterstellung und die Definition der Soll-Prozesse tragen. \cite{26}\cite{27}
    \item \uline{Externe Dienstleister} können als Berater oder Projektsteuerer ergänzend hinzugezogen werden, falls intern nicht ausreichend Know-how oder Kapazitäten verfügbar sind. \cite{28}
\end{itemize}
Die wichtigsten \uline{Ziele} des Migrationsprojekts lassen sich wie folgt beschreiben:
\begin{itemize}
    \item Eine klare \uline{Zieldefinition} ist zu Projektbeginn unerlässlich, um während der Laufzeit Orientierung zu bieten und nach Abschluss eine Erfolgskontrolle zu ermöglichen. \cite{29}\cite{30}
    \item Zu den primären \uline{organisatorischen Zielen} gehören die Verkürzung von Durchlaufzeiten, die Integration manueller Arbeitsabläufe in das System und die Neugestaltung von Prozessen auf Basis von \uline{Best Practices}. \cite{31}\cite{32}
    \item \uline{Wirtschaftliche Ziele} konzentrieren sich auf die Erreichung eines positiven \uline{Return on Investment (\textit{\acs{RoI}})} durch signifikante Kosteneinsparungen (\textit{\acs{z. B.} Verringerung des Materialaufwands}) und Produktivitätssteigerungen. \cite{33}\cite{34}\cite{35}\cite{36}\cite{37}
    \item \uline{Technische Ziele} beinhalten die Ablösung veralteter Software durch moderne, skalierbare Architekturen sowie die Gewährleistung einer hohen Datenaktualität durch die Erfassung von Vorgängen in \uline{Echtzeit}. \cite{32}\cite{38}\cite{39}\cite{40}
    \item Ein wesentliches Ziel ist zudem die Schaffung von \uline{Transparenz} über alle Unternehmensressourcen und Werteflüsse hinweg, um eine fundierte Basis für das Controlling und Managemententscheidungen bereitzustellen. \cite{31}\cite{41}
    \item Zuletzt soll die langfristige \uline{Wandlungsfähigkeit} sichergestellt werden, damit das \acs{ERP}-System auch künftige Marktveränderungen und organisatorische Umbrüche effizient abbilden kann. \cite{2}\cite{4}
\end{itemize}
\section{Was sind die fünf Unterschiede und Gemeinsamkeiten in den Funktionen eines \acs{SCM}- und eines \acs{ERP}-Systems?} \par Hier sind fünf wesentliche Unterschiede und Gemeinsamkeiten in den Funktionen und Merkmalen von \acs{ERP}- und \acs{SCM}-Systemen aufgeführt: \par \textbf{Unterschiede}
\begin{enumerate}[label=\arabic*.), leftmargin=*]
    \item \textbf{Fokus der Betrachtung:} Während \acs{ERP}-Systeme primär auf die \uline{unternehmensinterne Integration} und die Verwaltung interner Ressourcen ausgerichtet sind, fokussieren \acs{SCM}-Systeme auf die \uline{gesamte Lieferkette} und unternehmensübergreifende Wertschöpfungsnetzwerke. \cite{42}
    \item \textbf{Art der Planung:} Die Planung in \acs{ERP}-Systemen erfolgt meist \uline{sukzessiv} \acs{bzw.} \uline{sequenziell} nach dem \acs{MRP} II-Konzept, wohingegen \acs{SCM}-Systeme eine \uline{simultane Planung} (\textit{Advanced Planning and Scheduling – \acs{APS}}) ermöglichen. \cite{42}
    \item \textbf{Aktualität und Frequenz:} \acs{ERP}-Systeme arbeiten typischerweise mit täglichen oder wöchentlichen Intervallen, während \acs{SCM}-Systeme eine \uline{ständige Echtzeitverarbeitung} der Informationen erfordern. \cite{42}
    \item \textbf{Planungsgeschwindigkeit:} Ein Planungsvorgang in einem \acs{ERP}-System kann mehrere Stunden dauern, während \acs{SCM}-Systeme darauf ausgelegt sind, Planungsläufe in \uline{Minuten} abzuschließen. \cite{42}
    \item \textbf{Zielsetzung:} \acs{ERP}-Systeme dienen als Rückgrat der Informationsverarbeitung zur Administration interner Abläufe, während \acs{SCM}-Systeme spezifisch darauf zielen, netzwerkweite Probleme wie den \uline{Bullwhip-Effekt} durch Transparenz zu minimieren. \cite{42}\cite{43}
\end{enumerate}
\textbf{Gemeinsamkeiten}
\begin{enumerate}[label=\arabic*.), leftmargin=*]
    \item \textbf{Ressourcenmanagement:} Beide Systeme befassen sich grundlegend mit der \uline{Verwaltung} und \uline{Optimierung von Ressourcen}, wobei insbesondere der Materialfluss eine zentrale Rolle spielt. \cite{44}\cite{45}\cite{46}
    \item \textbf{Modularer Aufbau:} Sowohl \acs{ERP}- als auch \acs{SCM}-Systeme weisen eine \uline{modulare Architektur} auf, die es Unternehmen ermöglicht, einzelne Funktionsbereiche je nach Bedarf einzusetzen. \cite{47}\cite{48}
    \item \textbf{Prozessintegration:} Beide Systemkategorien unterstützen die \uline{horizontale Integration} von Geschäftsprozessen, um den Informationsfluss entlang der Wertschöpfungskette zu verbessern. \cite{49}\cite{50}\cite{51}\cite{52}
    \item \textbf{Entscheidungsunterstützung:} Sie dienen als \uline{strategische Informationswerkzeuge}, die dem Management auf Basis von Datenanalysen und Berichten fundierte Entscheidungen ermöglichen. \cite{53}\cite{54}
    \item \textbf{Transparenz und Kostensenkung:} Ein wesentliches gemeinsames Ziel ist die \uline{Schaffung von Transparenz}, um Ineffizienzen (\textit{\acs{z. B.} Überbestände}) aufzudecken und Kosten nachhaltig zu senken. \cite{55}\cite{56}\cite{57}
    \end{enumerate}

% ==============================================================================
% =================================== Aufgabe 1 =================================
% ==============================================================================